\section{Research target and the new exact families}
\label{clc:sec:scope}

The ProveIt manuscript already contains a substantial theory of
polylogarithmic identities, generalized harmonic sums, Hurwitz order
jets, Stieltjes distributions, and Gamma correlations \cite{Repo}.
The preceding \emph{Reciprocal Lerch Convolutions} report establishes
entire-order addition, finite central-factorial formulas, and all
logarithmic moments for equal logarithmic scales \cite{RLC}. Its proposed
research questions explicitly ask for commensurate dilations and for a
normally convergent unequal-shift expansion with independent order jets.
These are our principal targets. Re-proving an equal-scale identity is not
counted as a new result.

The central distinction is between a spectral order and a logarithmic
scale. If the argument $y$ is replaced by $y^q$, the bilateral transform
changes its cosecant period and also introduces a factor $q^s$. That
factor cannot be ignored. Once it is removed, the new periods can be
combined into a finite pole lattice exactly when their ratios are rational.
A separate expansion handles unequal affine centers without pretending
that a finite integer-order partial fraction identity can be differentiated
in its integer index.

\subsection{Two accessible consequences}
For arbitrary complex $s,t$, write $W=s+t$. One result proved below is
\begin{equation}\label{clc:eq:intro-beta}
 \boxed{\displaystyle
 \int_0^\infty \Li_s(-y)\Li_t(-y^{-2})\frac{\dd y}{y}
 =2^{t-1}\pi\beta(W)-2^{-s-1}W\eta(W+1).}
\end{equation}
Here $\beta$ is the Dirichlet beta function and $\eta$ the Dirichlet eta
function. Both sides are entire in the two orders, and the integral is
ordinary and absolutely convergent. In particular, with
$G=\beta(2)$ denoting Catalan's constant,
\begin{equation}\label{clc:eq:intro-log}
 \int_0^\infty \log(1+y)\log(1+y^{-2})\frac{\dd y}{y}
 =\pi G-\frac38\zeta(3).
\end{equation}
The two orders in \eqref{clc:eq:intro-beta} do not enter symmetrically.
The normalization restoring pure order addition is multiplication of the
integral by $2^{-t}$.

A different consequence does not require rationally related scales.
For any positive real $q_1,q_2,q_3$,
\begin{equation}\label{clc:eq:intro-cubic}
 \boxed{\displaystyle
 \int_{\R^2}
 \frac{1}{1+e^{-q_1x}}\,
 \frac{1}{1+e^{-q_2y}}\,
 \frac{1}{1+e^{q_3(x+y)}}\dd x\dd y
 =\frac{\pi^2}{12}\left(q_1^{-2}+q_2^{-2}+q_3^{-2}\right).}
\end{equation}
This is one case of an all-depth negative-integer parity theorem. It
illustrates that the absence of a common pole lattice does not rule out
exceptional exact central evaluations.

\subsection{Conventions}
We use the standard Hurwitz and Lerch functions \cite{DLMF25}:
\[
 \zeta(v,a)=\sum_{n=0}^\infty(n+a)^{-v},\qquad
 \Phi(z,v,a)=\sum_{n=0}^\infty\frac{z^n}{(n+a)^v},
\]
initially in their convergence domains and then by analytic continuation.
Throughout, $\Rea a>0$ unless another domain is explicitly stated; the
powers use the logarithm on the right half-plane. The rising factorial is
$(v)_j=v(v+1)\cdots(v+j-1)$, with $(v)_0=1$. A superscript $[h]$ means
an order derivative. A superscript $(h)$ on $\psi$ means an ordinary
argument derivative of the digamma function.

Our generalized Stieltjes convention is
\begin{equation}\label{clc:eq:stieltjes-convention}
 \zeta(1+v,a)=\frac1v+
 \sum_{n=0}^\infty\frac{(-1)^n}{n!}\gamma_n(a)v^n,
 \qquad \gamma_0(a)=-\psi(a).
\end{equation}
We repeatedly use the entire functions
\begin{align}
 \EE_a(v)&=v\zeta(v+1,a),&\EE_a(0)&=1,\label{clc:eq:E}\\
 \eta(v,a)&=\Phi(-1,v,a)
 =2^{-v}\left[\zeta\left(v,\frac a2\right)
                  -\zeta\left(v,\frac{a+1}2\right)\right].
 \label{clc:eq:eta}
\end{align}
In particular $\eta(v)=\eta(v,1)$, while
$\beta(v)=4^{-v}[\zeta(v,1/4)-\zeta(v,3/4)]$.
The singularities in these formulas are filled in before differentiation.

\subsection{Relationship to classical methods and remaining conjectures}
Euler's beta integral, Gamma reflection, convolution, pole summation,
and Bell-polynomial differentiation are classical tools; the use of real
convolution for Fermi--Dirac-type integrals is explicitly represented in
Selvaggi and Selvaggi \cite{SS}. We do not claim to have invented these
methods. The contribution here is the all-complex-order commensurate-scale
closure, its precise hypotheses, pole-free normalization, and the derived
families of identities and independent-center continuations.

The inspected canonical README records $S_4$ as proved and $S_6$ and the
revised $S_8$ candidate as conjectural. Those statuses are unchanged here.
A small residual is not an equality proof. The two research questions we
answer are the analytic continuation problems just described, not those
specific arithmetic reduction conjectures. Complete literature priority
for every special value has not been established.
